\subsection{Data and training}
\label{app:benchmark_methods}
\label{app:benchmark_protocols}
Tables~\ref{tab:benchmark_protocols} and~\ref{tab:benchmark_search} give the data and training settings. The repository contains the preprocessing and training code, together with the recovered settings and data splits. Missing settings are identified below.

\begin{table}[ht]
\centering
\caption{Data and training settings. M and T denote MLP and Transformer. $E$ is the epoch count and WD is weight decay. All rows use batch size 128. Counts are training/validation/test examples. The standalone FI-2010 counts follow the available segment manifest; its match to the original trial manifest is unverified.}
\label{tab:benchmark_protocols}
\small
\setlength{\tabcolsep}{4pt}
\renewcommand{\arraystretch}{1.2}
\begin{tabular}{@{}p{0.18\textwidth}p{0.37\textwidth}p{0.20\textwidth}p{0.18\textwidth}@{}}
\toprule
Dataset & Input and split & Counts & Models; $E$; WD \\
\midrule
Speech Commands & 35 classes; $64\!\times\!64$ log-mel image; balanced random clips from official training subset & 20,020/5,005/10,010 & M,T; 50; $0$ \\
Tiny ImageNet 20k & 200 classes; $3\!\times\!64\!\times\!64$ RGB; balanced subsets of official training images & 20,000/5,000/10,000 & M,T; 75; $0$ \\
Tiny ImageNet 80k & Same image pool; separately drawn balanced split & 80,000/10,000/10,000 & M,T; 75; $0$ \\
FI-2010 benchmark & 3 classes; $100\!\times\!40$ book window; one stock, forward split with 100-window gaps & 40,000/10,000/20,000 & M; 50; $0$ \\
FI-2010 standalone & $128\!\times\!40$ windows within each stock/day; train day 6, validation day 7, test days 8--10 & 38,467/36,661/137,532 & M; 100; $0$ \\
\bottomrule
\end{tabular}
\end{table}

\paragraph{Preparing the inputs.}
Continuous inputs use coordinate-wise training-set means and population standard deviations; a standard deviation below $10^{-8}$ is replaced by one. There is no data augmentation. Speech waveforms use the first channel, resample to 16\,kHz when needed, and truncate or zero-pad to one second. The mel transform uses FFT size 1024, hop 250, 64 mel bands, and an 80\,dB range; the code retains the first 64 frames. A reconstruction of the random clip split finds that 4,979 of 5,005 validation clips and 9,967 of 10,010 test clips share a speaker with training. These results therefore measure performance on held-out clips, mostly from familiar speakers. The 80\,dB cutoff was applied across each 256-clip preprocessing batch. Tiny ImageNet uses RGB pixels divided by 255 and sorted class/file order. The 20k experiment computes normalization statistics in float32; the 80k experiment combines float64 moments in 64\,MiB chunks before casting the statistics to float32. Their holdout sets differ. Random class-balanced splits use NumPy's \texttt{default\_rng} with seed 20260812, sampling without replacement within each class.

The benchmark FI-2010 cache reads the first 40 feature rows of the no-auction Z-score release's ninth cumulative training fold, within stock ordinal 5 only; labels use zero-based row 148, the fifth supplied horizon, at the end of the window. Its sequential window-index ranges are $[0,40000)$, $[40100,50100)$, and $[50200,70200)$. The standalone study instead uses decimal-precision daily segments, all five stock ordinals, and label row 144 (the first horizon). Its 128-row windows never cross stock/day boundaries, and the last ten representations of each segment are purged. Its feature normalization uses day-6 rows only, excluding these tails, with float64 moments cast to float32.

\paragraph{Architectures.}
Benchmark MLPs have 12 hidden affine--ReLU--dropout blocks of width 256 and a separate affine classifier. Every affine layer, including the readout, has independent Gaussian weights of variance $1.98/\mathrm{fan\_in}$ and biases of variance $0.02$. Transformers have 12 pre-LayerNorm residual blocks, width 256, eight attention heads, and a ReLU feedforward sublayer of width 1024. The layer's dropout probability is applied separately to the attention output and feedforward output before residual addition; attention weights have no dropout. Learned absolute positions and a learned class token precede a final LayerNorm and class-token classifier. Speech and Tiny ImageNet use $8\!\times\!8$ patches (64 patches). Linear weights and position/class embeddings use truncated-normal initialization with standard deviation $0.02$; linear biases are zero. Patch convolutions retain the library's default initialization. Standalone MLPs use the same Gaussian affine initialization with depths $\{6,12\}$ and widths $\{256,512\}$.

\paragraph{Training and evaluation.}
Benchmarks use AdamW, zero weight decay, unweighted cross-entropy, and no warmup. At epoch $e\in\{0,\ldots,E-1\}$ the learning rate is
\[
\eta_e=\eta_0\left[f+\frac{1-f}{2}\left(1+\cos\frac{\pi e}{E-1}\right)\right],
\qquad f=10^{-3}.
\]
Transformers clip the gradient norm at one; benchmark MLPs do not clip. Standalone MLPs use Adam, norm clipping at one, and the same cosine formula with $f=10^{-2}$. Both optimizers retain $(\beta_1,\beta_2)=(0.9,0.999)$ and $\epsilon=10^{-8}$. Training shuffles examples and retains the last partial batch. Benchmark CUDA training uses autocast (bfloat16 when supported, otherwise float16); standalone training uses float32 with deterministic algorithms and TF32 disabled.

We evaluate each run on test data using the weights from the epoch with the lowest validation loss, taking the earlier epoch in a tie. The first five seeds include final-epoch test results only for Tiny ImageNet 80k; the five additional seeds include them for every dataset. Final results in Table~\ref{tab:frontloaded_results} therefore use ten pairs for Tiny ImageNet 80k and five elsewhere. Test performance was not recorded at every epoch.

\par\medskip
\noindent\begin{minipage}{\textwidth}
\begingroup
\centering
\captionof{table}{Hyperparameter searches and seeds. The grids and selection rules come from the training code, with settings for the benchmark runs checked against saved configuration hashes. Additional seeds keep each profile's settings unchanged.}
\label{tab:benchmark_search}
\small
\setlength{\tabcolsep}{4pt}
\renewcommand{\arraystretch}{1.2}
\begin{tabular}{@{}p{0.19\textwidth}p{0.49\textwidth}p{0.26\textwidth}@{}}
\toprule
Experiment & Search and chosen settings & Seeds \\
\midrule
Zero weight decay & Per-profile LR search at $\bar p=0.10$: MLP $\{3\!\times\!10^{-5},10^{-4},3\!\times\!10^{-4},10^{-3},3\!\times\!10^{-3}\}$; Transformer replaces the first value by $5\!\times\!10^{-5}$. Then search $\bar p\in\{0.05,0.10,0.15,0.20\}$ at that profile's selected LR. No dropout has its own LR search. & LR: $0$; budget: $0,1,2$; all profiles: $100$--$104$; extension: $105$--$109$ \\
FI-2010 standalone & Uniform-only 50-epoch LR screen over $\{3\!\times\!10^{-5},10^{-4},3\!\times\!10^{-4},10^{-3}\}$; keep one LR per depth and width across seven profiles, with $\bar p=0.10$ except no dropout. & LR: $0,1,2$; 100-epoch confirmation: $100$--$104$ \\
\bottomrule
\end{tabular}
\par
\endgroup

\paragraph{Choosing the settings.}
For each setting, we average the lowest validation cross-entropy reached by each seed and choose the setting with the lowest average. Ties favor the lower learning rate, then the lower budget in the budget sweep. These hyperparameters are chosen before running the confirmation seeds. For Table~\ref{tab:frontloaded_results}, we choose the front-loaded profile by the same average over the five confirmation seeds, considering only profiles with all five paired runs complete. This profile choice is retrospective. Additional seeds keep that profile and its hyperparameters fixed, and the same profile is used for test results at the epoch with the lowest validation loss and at the final epoch.

\paragraph{Uncertainty across seeds.}
We pair profiles by seed, using separate random streams for initialization, minibatch order and dropout. The nominal 95\% intervals use Fieller's method for relative loss reductions and Student-$t$ intervals for accuracy gains. They measure variation across seeds for the chosen split and settings, without accounting for the earlier selection or multiple comparisons. Both use $n-1$ degrees of freedom for $n$ independent seed pairs. Fieller intervals assume approximately bivariate-normal paired losses; Student-$t$ intervals assume approximately normal paired accuracy differences.

\paragraph{Missing settings.}
The complete search results could not be recovered. For standalone FI-2010, the selected learning rates and hashes of the original runs, source and data are also missing. The released code repeats the documented search and training procedure, with missing values marked in the repository. Source hashes identify the recovered benchmark code; results need not be bitwise identical across library versions or devices.

\paragraph{Additional seeds.}
Seven benchmark comparisons add seeds $105$--$109$ to the original $100$--$104$, giving ten paired test measurements at the epoch with the lowest validation loss for each row in Table~\ref{tab:frontloaded_results}. The extension also brings the original CIFAR-10 ReLU comparison from three to ten pairs (new seeds $45$--$51$) and the both-block ViT comparison from five to ten (new seeds $47$--$51$). These 94 additional runs use the same profiles, hyperparameters and data splits, without a new search. We checked the benchmark settings and hashes of the full cached data splits before combining results. The additional runs used H100 GPUs, while the CIFAR experiments reported A100s. The ReLU implementation also required a change to tensor storage before flattening, preserving the order of the features. The repository contains the settings and measurements, with the initial and additional seeds listed separately.
\end{minipage}
\par\medskip
